\documentclass{article}
\usepackage[utf8]{inputenc}
\usepackage[italian]{babel}
\usepackage{xcolor}
\usepackage[ruled,lined,linesnumbered]{algorithm2e}

\title{Alberi}
\author{Riccardo Di Michele}

\begin{document}
\maketitle

\section*{Introduzione}

\section*{Rappresenzationi per alberi}

\subsection*{Vettore posizionale}
Questa tipologia è perfetta per gli alberi \textbf{d-ari},
ovvero quelle strutture i cui figli di un nodo sono sempre una quantità
$d$, eccetto per i nodi dell'ultimo livello. 
Preso un albero $T(n, d)$ d-ario, con nodi numerati da $0$ a $n - 1$, un 
\textbf{vettore posizionale} $P$ è un array di dimensioni $n$ dove
l'elemento $P[v]$ conserva gli elementi del corrispondente nodo $v$. \\
L'ordine in cui questi elementi sono conservati all'interno del vettore è 
molto importante poichè grazie al loro ordine possiamo facilmente trovare
qualsiasi figlio del nodo $v_i$ tramite questa formula:
\[ P[d \cdot v + i] \]

dove, ricordiamo, $d$ essere il numero dei figli di ogni nodo, $v$ il nodo 
di cui vogliamo trovare i figli, e $i$ sarà un valore compreso tra $1$ e $d$.
\\
\\
Per esempio in un albero ternario ($d = 3$) se volessi trovare i figli del 
nodo $v = 4$ allora avrei:
\begin{itemize}
  \item $P[3 \cdot 4 + 1] \rightarrow P[13]$, come indice del primo figlio 
  \item $P[3 \cdot 4 + 2] \rightarrow P[14]$, come indice del secondo figlio
  \item $P[3 \cdot 4 + 3] \rightarrow P[15]$, come indice del terzo figlio
\end{itemize}

%TODO esempio albero%

\end{document}